% Part: first-order-logic
% Chapter: axiomatic-proofs
% Section: identity

\documentclass[../../../include/open-logic-section]{subfiles}

\begin{document}

\olfileid{fol}{axd}{ide}

\olsection{带等词的推导}

为了在推导中容纳 $\eq$，我们只需增加新的公理模式。
推导及 $\Proves$ 的定义保持不变，只是额外允许使用这些新公理。

\begin{defn}[等词公理]
\begin{align}
\ollabel{ax:id1} & \eq[t][t], \\
\ollabel{ax:id2} & \eq[t_1][t_2] \lif (!B(t_1) \lif
  !B(t_2)),
\end{align}，其中 $t$, $t_1$, $t_2$ 为任意闭项。
\end{defn}

\begin{prop}\ollabel{prop:sound}
  公理 \olref{ax:id1} 和 \olref{ax:id2} 是有效的。
\end{prop}

\begin{proof}
练习。
\end{proof}

\begin{prob}
证明 \olref[fol][axd][ide]{prop:sound}。
\end{prob}

\begin{prop}\ollabel{prop:iden1}
 对任意项 $t$ 和集合~$\Gamma$，$\Gamma \Proves \eq[t][t]$。
\end{prop}

\begin{prop}\ollabel{prop:iden2}
  若 $\Gamma \Proves !A(t_1)$ 且 $\Gamma \Proves
  \eq[t_1][t_2]$，则 $\Gamma \Proves !A(t_2)$。
\end{prop}

\begin{proof}
公式
\[
(\eq[t_1][t_2] \lif (!A(t_1) \lif !A(t_2)))
\]
是公理~\olref{ax:id2} 的一个实例。结论由两次应用~\MP 推出。
\end{proof}

\end{document}
